\documentclass[11pt]{article}
\usepackage[margin=1in]{geometry}
\usepackage{amsmath,amssymb,amsthm}
\usepackage[T1]{fontenc}
\usepackage{lmodern}
\usepackage{microtype}
\usepackage[colorlinks=true,allcolors=blue]{hyperref}
\hypersetup{pdftitle={A coefficient normalization for positive definite forms},
pdfauthor={Alec Kriebel},pdfsubject={Brandes's Oberwolfach Problem 11},
pdfkeywords={positive definite forms, polarization, coefficient normalization}}
\newtheorem{theorem}{Theorem}
\newtheorem{lemma}{Lemma}
\newtheorem{corollary}{Corollary}
\newcommand{\R}{\mathbb R}
\newcommand{\Q}{\mathbb Q}
\newcommand{\GL}{\mathrm{GL}}
\title{A coefficient normalization for positive definite forms}
\author{Alec Kriebel\\[3pt]\small Independent researcher\\
\small\href{https://orcid.org/0009-0001-9320-500X}{ORCID: 0009-0001-9320-500X}}
\date{23 September 2026 \quad (version 1.0)}
\begin{document}
\maketitle
\begin{abstract}
Every positive definite real homogeneous form admits linear coordinates in
which its symmetric ordered coefficients have diagonal entries equal to one
and all mixed entries strictly between zero and one. This answers Brandes's
question in Problem 11 of the 2019 Oberwolfach report on analytic number
theory, including its absolute-value inequality. The proof chooses a basis
close to a direction minimizing the form on a Euclidean sphere. The desired
coefficient inequalities then follow from a positive quadratic term in a
polynomial difference.
\end{abstract}

\section{Statement and coefficient convention}
Let $p\colon\R^m\to\R$ be a homogeneous polynomial of positive degree $d$, with
$p(x)>0$ for $x\ne0$. Then $d$ is even and $d\ge2$. Let $A$ be the unique
symmetric $d$-linear form satisfying $p(x)=A(x,\ldots,x)$. Throughout, a
coefficient means a \emph{symmetric ordered coefficient}:
\[
 p(t)=\sum_{j_1,\ldots,j_d=1}^m
 A(e_{j_1},\ldots,e_{j_d})t_{j_1}\cdots t_{j_d}.
\]
Thus the ordinary coefficient of $t^\alpha$ is
$\binom{d}{\alpha}A(e_1^{\alpha_1},\ldots,e_m^{\alpha_m})$.
Repeated arguments are denoted by superscripts inside $A$.

Brandes asked whether a real invertible linear change of variables can make
each coefficient's absolute value at most the geometric mean of its
corresponding diagonal entries \cite[p.~3182, Problem~11]{owr}.
The earlier discussion of pseudo-diagonal forms in \cite{brandes} includes
the quadratic case in Lemma~2.1. We prove the general basis-existence assertion
with the ordered coefficient convention above.

\begin{theorem}\label{thm:main}
Let $m\ge1$ and let $p$ be a positive definite real homogeneous form of
degree $d\ge2$. There is a basis $z_1,\ldots,z_m$ of $\R^m$ such that
\[
 p(z_i)=1\quad(1\le i\le m),\qquad
 0<A(z_{j_1},\ldots,z_{j_d})<1
\]
whenever the indices $j_1,\ldots,j_d$ are not all equal. In particular, in
these coordinates the requested absolute-value bound holds for every tuple,
with equality precisely for the pure tuples.
\end{theorem}

The change of variables is in $\GL_m(\R)$; no orthogonality or quantitative
conditioning constraint is imposed. For $m=1$ the assertion about mixed
tuples is vacuous.

\section{Proof}
Fix an auxiliary Euclidean norm. Choose a unit vector $v$ where $p$ attains
its sphere minimum $c>0$, and first replace $p,A$ by $p/c,A/c$.
The coefficient inequalities are invariant under this positive common
scaling. Hence for the intermediate construction we have
\begin{equation}\label{eq:comparison}
 p(v)=1,\qquad p(x)\ge \lVert x\rVert^d\quad(x\in\R^m).
\end{equation}
Write $Q(a,b)=A(v^{d-2},a,b)$.

\begin{lemma}\label{lem:tangent}
For every $u\perp v$,
\[
 A(v^{d-1},u)=0,\qquad Q(u,u)\ge\frac{\lVert u\rVert^2}{d-1}.
\]
\end{lemma}
\begin{proof}
By \eqref{eq:comparison}, the function
$p(v+tu)-(1+t^2\lVert u\rVert^2)^{d/2}$ has a minimum of zero at $t=0$.
Its first derivative is $dA(v^{d-1},u)$, and its second derivative is
$d(d-1)Q(u,u)-d\lVert u\rVert^2$. They are respectively zero and
nonnegative, as claimed.
\end{proof}

Choose an orthonormal basis $w_1,\ldots,w_{m-1}$ of $v^\perp$, set $w_m=0$,
and define
\begin{equation}\label{eq:basis}
 x_i(\varepsilon)=v+\varepsilon w_i\qquad(1\le i\le m).
\end{equation}
For every $\varepsilon\ne0$ these vectors are independent: the tangential
component of $\sum_i a_i x_i=0$ forces $a_1=\cdots=a_{m-1}=0$, and then
the radial component forces $a_m=0$.

For a fixed ordered tuple $J=(j_1,\ldots,j_d)$ put $a_r=w_{j_r}$ and
\[
 N_J(\varepsilon)=A(v+\varepsilon a_1,\ldots,v+\varepsilon a_d),\qquad
 F_J(\varepsilon)=\prod_{r=1}^d p(v+\varepsilon a_r)-N_J(\varepsilon)^d.
\]
Both expressions are polynomials in $\varepsilon$. Lemma~\ref{lem:tangent}
and multilinearity give
\begin{align*}
 N_J(\varepsilon)
 &=1+\varepsilon^2\sum_{r<s}Q(a_r,a_s)+O(\varepsilon^3),\\
 p(v+\varepsilon a_r)
 &=1+\binom d2\varepsilon^2 Q(a_r,a_r)+O(\varepsilon^3).
\end{align*}
Consequently
\begin{align}
 F_J(\varepsilon)
 &=\varepsilon^2\left(\frac{d(d-1)}2\sum_rQ(a_r,a_r)
              -d\sum_{r<s}Q(a_r,a_s)\right)+O(\varepsilon^3)\notag\\
 &=\frac d2\varepsilon^2\sum_{r<s}Q(a_r-a_s,a_r-a_s)
       +O(\varepsilon^3).\label{eq:gap}
\end{align}
The last equality follows by expanding each pair difference; each diagonal
term occurs in $d-1$ pairs.

If $J$ is mixed, some $a_r-a_s$ is a nonzero vector in $v^\perp$.
Lemma~\ref{lem:tangent} makes the coefficient of $\varepsilon^2$ in
\eqref{eq:gap} strictly positive. Thus $F_J(\varepsilon)>0$ for all
sufficiently small positive $\varepsilon$. There are only finitely many
tuples, so one positive $\varepsilon$ works for all mixed tuples.
For a pure tuple, $F_J$ is identically zero. Also $N_J(0)=1$ for every
tuple; by continuity and finiteness, decreasing $\varepsilon$ if needed
ensures $N_J(\varepsilon)>0$ for all tuples.

Since $d$ is even, these conclusions imply
\[
 0<A(x_{j_1},\ldots,x_{j_d})
 \le \prod_{r=1}^d p(x_{j_r})^{1/d},
\]
strictly for mixed tuples. Restore the original $p$ and $A$: both sides
scale by $c$, so the inequality is unchanged. Finally set
$z_i=x_i/p(x_i)^{1/d}$ using the original $p$. Positive column rescaling
preserves independence, makes $p(z_i)=1$, and proves
Theorem~\ref{thm:main}. The matrix $S$ with columns $z_i$ has
$p'(t)=p(St)$ and polarization
$A'(e_{j_1},\ldots,e_{j_d})=A(z_{j_1},\ldots,z_{j_d})$, as required.
\hfill$\square$

\section{Scope and verification}
The basis in \eqref{eq:basis} may be very close to singular. This is allowed
by the question. Strict inequalities and nonzero determinant are open
conditions, so the unnormalized basis $x_i$ can also be chosen rational
by density of $\Q^{m\times m}$ in $\R^{m\times m}$. The final unit-diagonal
normalization need not preserve rationality. No claim is made here for
merely nonnegative forms or for unimodular integral changes of variables.

The conclusion concerns the existence of suitable coordinates. For example,
$p(x,y)=x^4+12x^2y^2+y^4$ is positive definite but has
$A(e_1,e_1,e_2,e_2)=2>1$ in the original coordinates. The multinomial
factor cannot be dropped when translating to ordinary monomial coefficients.

The accompanying open package contains an independent proof using
positive definite bilinear slices and Cauchy--Schwarz, internal adversarial
reviews, a dated priority audit, and a Python standard-library verifier.
The verifier checks 12 rational examples exactly, including a nonconvex
positive definite quartic, and cross-checks coefficients by polarization.
These finite checks supplement the analytic proof; they are not a formal
verification of the universal theorem. No separate computational assumption
is needed in the proof.

\paragraph{Provenance and status.}
This is an unrefereed preprint. Generative AI supplied the initial candidate
argument and assisted the internal independent checks, literature search,
and writing. A bounded public-literature search completed on 23 September
2026 found no earlier general resolution; it does not establish first
priority. The two direct antecedents are cited below.
Paper, source, and verification materials are available at
\href{https://aleckriebel.github.io/Math/papers/brandes-coefficient-normalization/}{the project page}.

\begin{thebibliography}{9}
\bibitem{brandes}
J.~Brandes, \emph{Forms representing forms: the definite case},
J. London Math. Soc. (2) \textbf{92} (2015), no.~2, 393--410.
\href{https://doi.org/10.1112/jlms/jdv028}{doi:10.1112/jlms/jdv028}.
\bibitem{owr}
H.~L. Montgomery, \emph{Problem Session}, in \emph{Analytic Number Theory},
Oberwolfach Reports \textbf{16} (2019), no.~4, Report~50, pp.~3179--3183;
Problem~11 (Brandes), p.~3182.
\href{https://doi.org/10.4171/OWR/2019/50}{doi:10.4171/OWR/2019/50}.
\end{thebibliography}
\end{document}
